\documentclass[12pt,a4paper]{article}
\usepackage[UTF8,heading=true]{ctex}
\usepackage{geometry}
\usepackage{amsmath,amssymb}
\usepackage{array}
\usepackage{enumitem}
\usepackage{fancyhdr}
\usepackage{setspace}
\usepackage{hyperref}

\geometry{left=2.5cm,right=2.5cm,top=2.5cm,bottom=2.5cm}
\setlength{\headheight}{15pt}
\setstretch{1.5}
\raggedbottom
\setlength{\emergencystretch}{3em}

\pagestyle{fancy}
\fancyhf{}
\fancyhead[C]{权利要求书（审查稳版）}
\fancyfoot[C]{\thepage}

\begin{document}

\begin{center}
    {\Huge\heiti 权利要求书（审查稳版）}
\end{center}

\vspace{1em}

\noindent\textbf{发明名称：}一种面向 6G 网络切片的超低时延抗量子自适应握手方法、系统、设备及存储介质

\vspace{1em}

\textbf{权利要求1.}一种面向 6G 网络切片的抗量子自适应握手方法，其特征在于，包括：

\begin{enumerate}[label=\arabic*., leftmargin=2em]
    \item 接收终端发送的握手请求消息，所述握手请求消息至少携带切片标识、终端能力信息以及终端支持的后量子密码算法集合；
    \item 根据所述切片标识对应的切片握手画像，并结合所述终端能力信息和链路状态信息，确定目标握手计划，所述切片握手画像至少限定是否允许早期业务、增强片段恢复阈值、增强片段总数以及增强阶段时限；
    \item 按照所述目标握手计划执行基础后量子密钥协商，获得基础份额，并基于所述基础份额导出基础会话密钥；
    \item 在所述切片握手画像允许的情况下，基于所述基础份额、切片标识和握手转录文本摘要导出早期业务密钥，并限定所述早期业务密钥仅用于切片策略预定义的白名单业务；
    \item 生成增强密钥材料，并对所述增强密钥材料执行全有或全无变换、秘密共享、纠删编码中的至少一种处理，以形成总数为 $n$ 且恢复阈值为 $t$ 的多个增强片段，其中 $t \leq n$；
    \item 在所述增强阶段时限内，通过后续握手报文、控制报文、首批业务报文或多路径子流中的至少一种承载路径发送或接收所述多个增强片段；
    \item 在接收到达到所述恢复阈值 $t$ 的有效增强片段时恢复所述增强密钥材料；
    \item 基于所述基础份额、所述增强密钥材料、所述握手转录文本摘要、所述切片标识和所述目标握手计划导出最终会话密钥；
    \item 将所述切片标识、握手画像标识、握手计划标识、算法包摘要、恢复阈值、增强片段总数、增强阶段时限以及早期业务范围绑定到所述握手转录文本中，以实现抗降级和抗跨切片重放保护。
\end{enumerate}

\vspace{0.5cm}

\textbf{权利要求2.}根据权利要求1所述的方法，其特征在于，所述切片握手画像还至少包括以下一项或多项：切片时延预算、最小安全等级、允许的算法集合、业务权限级别以及增强阶段失效窗口。

\vspace{0.5cm}

\textbf{权利要求3.}根据权利要求1所述的方法，其特征在于，所述基础后量子密钥协商所采用的参数集包括以下至少一种：单一后量子 KEM 参数集、多后量子 KEM 参数集、椭圆曲线密钥交换与后量子 KEM 的混合参数集。

\vspace{0.5cm}

\textbf{权利要求4.}根据权利要求1所述的方法，其特征在于，所述增强片段在发送前还附带片段序号、路径标识、计划标识或切片标识中的至少一种元数据，以便在多路径传输或重排序场景下进行恢复校验。

\vspace{0.5cm}

\textbf{权利要求5.}根据权利要求1所述的方法，其特征在于，所述白名单业务至少受切片类型、业务优先级和消息类型中的一项或多项限制，且所述白名单业务范围被写入所述握手转录文本的绑定内容。

\vspace{0.5cm}

\textbf{权利要求6.}根据权利要求1所述的方法，其特征在于，若在所述增强阶段时限内未达到所述恢复阈值 $t$，则根据切片策略执行以下至少一种操作：限制业务范围、降低会话权限、请求重发增强片段、发起补强重协商、撤销当前会话。

\vspace{0.5cm}

\textbf{权利要求7.}根据权利要求1所述的方法，其特征在于，所述绑定到握手转录文本中的内容还包括路径标识、片段顺序摘要和恢复票据摘要中的至少一种，以增强抗路径篡改能力。

\vspace{0.5cm}

\textbf{权利要求8.}根据权利要求1所述的方法，其特征在于，所述方法还包括：生成与所述切片标识、握手画像标识和策略摘要绑定的切片恢复票据，以用于同一切片中的后续快速重连，且所述切片恢复票据不能跨切片复用。

\vspace{0.5cm}

\textbf{权利要求9.}一种面向 6G 网络切片的抗量子自适应握手系统，其特征在于，包括：

\begin{enumerate}[label=\alph*), leftmargin=2em]
    \item 切片策略控制模块，用于存储不同网络切片对应的切片握手画像；
    \item 握手计划决策模块，用于根据切片标识、终端能力信息和链路状态信息选择目标握手计划；
    \item 基础份额协商模块，用于执行基础后量子密钥协商并导出基础会话密钥；
    \item 早期业务控制模块，用于在切片策略允许时导出早期业务密钥并限制其业务承载范围；
    \item 增强材料编码模块，用于生成增强密钥材料并形成满足阈值恢复条件的多个增强片段；
    \item 增强片段调度模块，用于在增强阶段时限内调度所述增强片段的发送、接收、重传或多路径分发；
    \item 增强材料恢复模块，用于在接收到达到恢复阈值的有效增强片段时恢复增强密钥材料；
    \item 最终密钥导出模块，用于导出最终会话密钥；
    \item 转录绑定保护模块，用于将切片相关参数、阈值参数和业务范围参数绑定到握手转录文本中。
\end{enumerate}

\vspace{0.5cm}

\textbf{权利要求10.}根据权利要求9所述的系统，其特征在于，所述增强材料编码模块被配置为对增强密钥材料先执行全有或全无变换，再执行纠删编码或秘密共享，以降低部分片段泄露时的信息暴露风险。

\vspace{0.5cm}

\textbf{权利要求11.}一种电子设备，其特征在于，包括处理器和存储器，所述存储器中存储有计算机程序，所述计算机程序在被所述处理器执行时，使得所述电子设备执行权利要求1至8任一项所述的方法。

\vspace{0.5cm}

\textbf{权利要求12.}一种计算机可读存储介质，其上存储有计算机程序，所述计算机程序在被处理器执行时，使得所述处理器执行权利要求1至8任一项所述的方法。

\end{document}
